\chapter{Authentication, sessions, and authorization}

Authentication proves who the user is; authorization decides what that user may do. After a successful sign-in, the session carries that identity across requests. Authentication is therefore not merely a login form. In a production web application, the identity backend, session lifecycle, second factor, authorization, CSRF protection, rate limiting, audit, and proxy trust form one system. RWLang separates these responsibilities: the \code{.rw} application declares route policy, while identity and session policy come from server configuration.

\section{The complete login path}

The built-in HTML login endpoints are:

\begin{lstlisting}
GET  /__rw/auth/login
POST /__rw/auth/login
POST /__rw/auth/logout
\end{lstlisting}

These are server endpoints, not routes defined by the application. The GET login page issues a CSRF token bound to the session. The POST login form expects exactly the \code{_csrf}, \code{username}, \code{password}, and \code{totp} fields.

At a high level, a successful login:

\begin{enumerate}
  \item validates form structure and the CSRF token;
  \item canonicalizes the username;
  \item applies the login rate limit;
  \item verifies local-auth or LDAP credentials;
  \item if required, verifies TOTP or, for local auth, a recovery code;
  \item rotates to a new authenticated session;
  \item writes an audit event;
  \item redirects to the root with \code{303 See Other}.
\end{enumerate}

Session rotation matters: after successful authentication, the anonymous session identifier is not kept. This is a basic defense against session fixation.

\section{Local auth: a separate identity database}

A single-node application may use the local-auth SQLite store. It is deliberately separate from the application database. RWLang typed SQL queries therefore cannot access credential tables.

For example, the secret file may be:

\begin{lstlisting}
/run/secrets/rwlang/local-auth-db-url
\end{lstlisting}

with contents:

\begin{lstlisting}
sqlite:///srv/rwlang/auth/users.db?mode=rwc
\end{lstlisting}

Initialize it with:

\begin{lstlisting}[language=bash]
/usr/local/bin/rwlang-cli auth init \
  --db-url-file /run/secrets/rwlang/local-auth-db-url
\end{lstlisting}

When creating a user, do not pass the password as a command-line argument; read it from a file:

\begin{lstlisting}[language=bash]
printf '%s\n' 'a-very-long-random-password' \
  > /run/secrets/rwlang/new-user-password
chmod 600 /run/secrets/rwlang/new-user-password

/usr/local/bin/rwlang-cli auth user-add \
  --db-url-file /run/secrets/rwlang/local-auth-db-url \
  --username alice \
  --password-file /run/secrets/rwlang/new-user-password \
  --role User \
  --role Editor
\end{lstlisting}

The local store uses Argon2id password hashing with a unique random salt. The current implementation accepts passwords between 12 and 1024 characters. Do not confuse this with application-side form validation: credential policy belongs to the identity backend.

\section{Account lifecycle through the operator CLI}

User state does not need to be modified with direct SQL. The CLI provides explicit operations:

\begin{lstlisting}[language=bash]
rwlang-cli auth password-set --db-url-file ... \
  --username alice --password-file ...

rwlang-cli auth disable   --db-url-file ... --username alice
rwlang-cli auth enable    --db-url-file ... --username alice
rwlang-cli auth roles-set --db-url-file ... --username alice \
  --role User --role Publisher
\end{lstlisting}

Disabling an account does not delete it. From the outside, failed login does not reveal whether the cause was a wrong password, nonexistent user, or disabled account, because such distinctions can become account-enumeration channels.

\section{Session invalidation: auth\_generation}

With local auth, credential and authorization changes increment the user's \code{auth_generation}. A session records the generation that was current when it was created.

On the next request, the server compares the session generation with the current value in the local-auth store. If they differ, the old session is invalidated and a fresh anonymous session is created.

Consequently, the following operator actions affect existing sessions too, not merely the next login:

\begin{itemize}
  \item password change;
  \item role change;
  \item account disable/enable;
  \item TOTP enrollment or disable.
\end{itemize}

This is particularly important for role revocation: there is no need to wait for the old session to expire naturally.

\section{Session store: memory or Redis}

An in-memory session store may be used for development and testing. In production, especially with multiple server instances, Redis should be the shared session backend.

\begin{lstlisting}
[redis]
url_file = "/run/secrets/rwlang/redis-url"
allow_insecure = false
\end{lstlisting}

Redis may hold several kinds of short-lived infrastructure state in RWLang:

\begin{itemize}
  \item shared sessions;
  \item TOTP replay protection;
  \item login rate limiting;
  \item public cache.
\end{itemize}

Redis is not a business source of truth. After disaster recovery, an empty session/cache state is acceptable: users sign in again and caches rebuild.

\section{Cookie policy}

In normal production mode the session cookie is named:

\begin{lstlisting}
__Host-rw_session
\end{lstlisting}

The server uses \code{Path=/}, \code{HttpOnly}, and \code{Secure}. The default is \code{SameSite=Lax}; with credentialed cross-origin configuration it becomes \code{SameSite=None; Secure}.

The explicit development escape hatch \code{--insecure-dev-cookies} changes the cookie name to \code{rw_session} and omits \code{Secure}. Never use this for an Internet-facing production deployment.

\section{Logout}

Logout is intentionally a POST operation:

\begin{lstlisting}[language=RWLang]
<form method="post" action="/__rw/auth/logout">
    <input type="hidden" name="_csrf" value="{{ csrfToken }}">
    <button type="submit">Logout</button>
</form>
\end{lstlisting}

The server verifies the CSRF token, invalidates the old session, creates a fresh anonymous session, sends a new cookie, and audits the logout event. Logout is therefore not simply a client-side cookie deletion.

\section{LDAP / LDAPS for enterprise identity}

For shared or enterprise identity, an LDAP backend may be used. Use an LDAPS URL in production, with service-bind credentials loaded from secret files.

Example configuration:

\begin{lstlisting}
[auth]
ldap_url = "ldaps://ldap.example.com:636"
ldap_search_base = "ou=people,dc=example,dc=com"
ldap_username_attribute = "uid"
ldap_service_bind_dn_file = "/run/secrets/rwlang/ldap-bind-dn"
ldap_service_bind_password_file = "/run/secrets/rwlang/ldap-bind-password"
require_totp = true
\end{lstlisting}

Do not configure local auth and LDAP as simultaneous primary identity backends in one process. That avoids ambiguous backend priority and username collisions.

With LDAP, role mapping should come from trusted server-side configuration. Client-supplied \code{role}, \code{username}, or similar fields are never authorization authority.

\section{TOTP and recovery codes}

Enroll TOTP for a local-auth user with:

\begin{lstlisting}[language=bash]
/usr/local/bin/rwlang-cli auth totp-enroll \
  --db-url-file /run/secrets/rwlang/local-auth-db-url \
  --username alice
\end{lstlisting}

The CLI prints the secret/URI data required for enrollment and the recovery codes. Recovery codes remain in the local-auth database only as hashes, and a successfully used code is consumed.

Disable TOTP with:

\begin{lstlisting}[language=bash]
/usr/local/bin/rwlang-cli auth totp-disable \
  --db-url-file /run/secrets/rwlang/local-auth-db-url \
  --username alice
\end{lstlisting}

If \code{require_totp = true}, a user with a correct password still cannot sign in without a usable second factor.

The login form's \code{totp} field may contain either a six-digit TOTP or, with local auth, a recovery code. With Redis, replay state is shared, so multiple server instances see the same TOTP replay protection.

\section{Route authorization}

Applications commonly declare three route-level policies:

\begin{lstlisting}[language=RWLang]
route account GET "/account" auth user => account;
route admin GET "/admin" auth role Admin => admin;
route security GET "/security" auth mfa => security;
\end{lstlisting}

Their meanings are:

\begin{itemize}
  \item \code{auth user}: any authenticated principal;
  \item \code{auth role Admin}: an authenticated principal with the trusted \code{Admin} role in the session;
  \item \code{auth mfa}: an authenticated principal that has successfully proven a second factor.
\end{itemize}

For an HTML route, an unauthenticated request may redirect to the login page. For a JSON API, the same condition appears as \code{401 Unauthorized} rather than an HTML redirect, which is more predictable for API clients.

\section{Using the principal on a page}

An authenticated page receives trusted session state:

\begin{lstlisting}[language=RWLang]
page fn account(ctx: PageContext) -> Result<Html, PageError> {
    return Ok(html {
        <main>
            <h1>Account</h1>
            <p>Signed in as {{ authPrincipal }}</p>
            <form method="post" action="/__rw/auth/logout">
                <input type="hidden" name="_csrf" value="{{ csrfToken }}">
                <button type="submit">Logout</button>
            </form>
        </main>
    });
}

route account GET "/account" auth user => account;
\end{lstlisting}

\code{authPrincipal}, session roles, and MFA state are server-side trusted state. Do not reconstruct them from hidden form fields or query parameters.

\section{A route role is not always enough}

A route-level role answers whether a user may enter an operation. Object-level authorization may require more.

For example, an editor may be allowed to modify only their own draft. In that case \code{auth role Editor} alone is insufficient: the action must compare the owner of the record loaded from the database with the trusted session principal.

The core rule is:

\begin{center}
\textbf{identity from the session, object state from the database, decision on the server}
\end{center}

Client-supplied \code{owner}, \code{actor}, or \code{role} fields are data at most, never authority.

\subsection*{The typed object-authorization guard}

The \code{authorize} guard may be used only on an already-loaded, \textbf{non-optional} model record. The compiler also verifies that the owner field exists and has type \code{String}. A \code{Result<Article?, DbError>} must therefore be handled explicitly first; authorization does not interpret a ``no record'' state.

\begin{lstlisting}[language=RWLang]
let article = loadArticle(db, id)?;
authorize article owner authorUsername or role Publisher;
\end{lstlisting}

A handler containing this guard cannot be exposed through a public route, and public cache cannot sit in front of object authorization. The owner or one of the listed trusted roles is allowed; every other authenticated actor receives \code{403 Forbidden}.

\section{Login rate limiting}

Login has its own rate-limit policy. For example:

\begin{lstlisting}
[auth]
login_max_attempts = 8
login_window_secs = 300
\end{lstlisting}

After trusted-proxy processing, the server uses the effective client IP together with the canonical username as part of the limiter key. This is why \code{trusted_proxy_cidrs} is also an authentication-security setting behind a reverse proxy.

A failed credential or second factor does not produce detailed failure information for the client. Use \code{429 Too Many Requests} for rate limiting and \code{503 Service Unavailable} when the authentication backend is unavailable.

\section{CSRF, Origin, and Fetch Metadata}

With session-cookie authentication, the browser sends the cookie automatically, so state-changing requests require CSRF protection. RWLang uses defense in depth:

\begin{itemize}
  \item a session-bound CSRF token;
  \item same-origin Origin/Referer policy;
  \item Fetch Metadata controls;
  \item strict request parsing;
  \item trusted-proxy policy.
\end{itemize}

CORS does not replace CSRF protection. If you allow credentialed CORS for a separate frontend origin, that is a distinct security policy and requires HTTPS.

\section{Authentication audit}

Login and logout events go to the audit log. Useful fields include request ID, actor, outcome, and effective client IP. The purpose of the audit log is event reconstruction, not credential debugging.

Never log:

\begin{itemize}
  \item passwords;
  \item TOTP secrets or recovery codes;
  \item session IDs or full Cookie headers;
  \item CSRF tokens;
  \item LDAP bind or DB/Redis credentials.
\end{itemize}

A login result may distinguish success, invalid credentials, invalid second factor, or rate limiting internally, while the client-facing message remains intentionally less detailed.

\section{Local auth or LDAP?}

\begin{center}
\begin{tabularx}{\textwidth}{@{}lYY@{}}
\toprule
 & Local auth & LDAP/LDAPS \\
\midrule
Identity store & separate SQLite DB & enterprise directory \\
Typical deployment & single-node & shared/enterprise identity \\
User administration & RWLang CLI & directory admin workflow \\
Role source & local-auth store & trusted mapping \\
TOTP & local enrollment + recovery & server-side TOTP policy/secrets \\
Multi-instance identity & not the target use case & yes \\
\bottomrule
\end{tabularx}
\end{center}

If the application needs public self-registration, email password reset, or self-service MFA enrollment, do not bolt a thin route onto the operator CLI. Those features require token lifecycle, email security, rate limiting, and their own recovery threat model.

\section{Rich text in an authenticated CMS}

An authenticated CMS editor interface still does not justify storing raw HTML. Editor authorization and content sanitization are separate defenses. Store Markdown as content and render it with the safe renderer:

\begin{lstlisting}[language=RWLang]
<article class="article-body">
    @markdown(article.body)
</article>
\end{lstlisting}

\code{@markdown} does not permit raw HTML and filters link targets through an allowlist.

\section{Production checklist}

Before enabling authentication in production, verify:

\begin{itemize}
  \item exactly one primary identity backend is active;
  \item the local-auth DB is separate from the application DB;
  \item passwords and service credentials come only from secret files;
  \item production sessions use Redis when multiple server instances run;
  \item cookies are Secure/HttpOnly and \code{--insecure-dev-cookies} is not enabled;
  \item \code{trusted_proxy_cidrs} covers only the real Apache/Nginx proxy addresses;
  \item the login limiter is configured;
  \item state-changing forms use CSRF protection;
  \item object-level authorization supplements route roles where necessary;
  \item credentials and session secrets never enter logs;
  \item account disable, role changes, and TOTP changes have been tested for session revocation.
\end{itemize}
